\Article{247}
\textsc{Theorema.}\quad\emph{Productum e binis periodis similibus independenter a numero $p$ componi potest per additionem periodorum similium et numerorum datorum.}

(Periodos similes vocamus, quae aeque multos terminos comprehendant sive ubi numerus $e$ est idem).

Exempl. Sit $n=7$, productum e periodis $(1)+(6)$ et $(2)+(5)$ erit (propter $(a)\times(b)=(a+b)$) $(3)+(6)+(8)+(11)$ sive constat e periodis $(3)+(4)$ et $(1)+(6)$.

\emph{Demonstr.}\quad Sit $\frac{\lambda}{e}=f$, atque periodi datae $(e\ast\alpha)$ et $(e\ast\beta)$ seu aggregata
\[
\begin{aligned}
&(\rho^\alpha)+(\rho^{\alpha+e})+(\rho^{\alpha+2e})+\ldots+(\rho^{\alpha+(f-1)e})\ \dotsfive\ P\\
&(\rho^\beta)+(\rho^{\beta+e})+(\rho^{\beta+2e})+\ldots+(\rho^{\beta+(f-1)e})\ \dotsfive\ Q
\end{aligned}
\]

Productum $PQ$ ex $f^2$ terminis constabit. Hi vero ita sunt ordinandi. Formentur $f$ series, quarum singulae ex $f$ terminis constent. Prima complectatur productum ipsius $P$ in $(\rho^\beta)$, secunda productum $P\pdot(\rho^{\beta+e})$ etc. etc. In prima serie primum locum occupet productum ex parte $(\rho^\alpha)$ oriundum, secundum productum ex $(\rho^{\alpha+e})$ et sic cetera deinceps; in secunda vero primus locus producto e parte $(\rho^{\alpha+e})$ oriundo tribuatur, secundus producto e parte $(\rho^{\alpha+2e})$ etc., ultimus denique producto e parte $(\rho^\alpha)$; tertia inchoet a producto e parte $(\rho^{\alpha+2e})$ et sic porro, post productum. e parte ultima sequatur productum e parte prima et secunda etc. etc., sive partibus successivis periodi $P$ per $1,2,3\ldots z$ et periodi $Q$ per $\mathrm{I},\mathrm{II},\mathrm{III},\dotsfour Z$ designati $PQ$ partes constituantur
\[
\begin{aligned}
&1\pdot\mathrm{I}+2\pdot\mathrm{I}+3\pdot\mathrm{I}+4\pdot\mathrm{I}+\dotstwo+z\pdot\mathrm{I}\\
&2\pdot\mathrm{II}+3\pdot\mathrm{II}+4\pdot\mathrm{II}+\dotsfour+1\pdot\mathrm{II}\\
&3\pdot\mathrm{III}+4\pdot\mathrm{III}+\ldots+1\pdot\mathrm{III}+2\pdot\mathrm{III}\\
&\qquad\text{etc.}\quad\text{etc.}
\end{aligned}
\]

Tunc omnes termini in singulis seriebus eundem locum occupantes in $f$ ordines colligantur; et dico
\SourceLeaf{216}
\par\noindent $1^\circ$ si aliquis terminus $\equiv1$, tum omnes ceteros eiusdem ordinis etiam fore $\equiv1$

\noindent $2^\circ$ quemvis ordinem, in quo nullus terminus $\equiv1$, periodum formare. --- Manifesto his demonstratis propositum consecuti erimus.

Forma generalis talis ordinis erit
\[
\begin{gathered}
(\rho^{\alpha+ke}+\rho^\beta),\quad(\rho^{\alpha+(k+1)e}+\rho^{\beta+e}),\quad(\rho^{\alpha+(k+2)e}+\rho^{\beta+2e}),\\
\ldots\ (\rho^{\alpha+(k+f-1)e}+\rho^{\beta+(f-1)e})
\end{gathered}
\]
potest enim pro $\rho^{\alpha+(k-1)e}$ etiam scribi $\rho^{\alpha+(k+f-1)e}$ propter $ef=\lambda$ et $\rho^\lambda\equiv1\;(\mathrm{mod.}\,n)$, et sic de antecedentibus. Ponatur $\rho^{\alpha+ke}+\rho^\beta\equiv\rho^\chi\;(\mathrm{mod.}\,n)$, quod est permissum, nisi forte $\rho^{\alpha+ke}+\rho^\beta$ per $n$ divisibilis\begingroup\renewcommand{\thefootnote}{*)}\footnote{Propositio paullo aliter exprimi debebit, si $n$ generaliter numeri primi potestatem denotat; quando vero est numerus primus, nihil immutandum.}\endgroup, poteritque ordo ita exhiberi $(\rho^\chi)$, $(\rho^{\chi+e})$, $(\rho^{\chi+2e})\ldots(\rho^{\chi+(f-1)e})$, qui manifesto est periodus $(e\ast\chi)$; si vero $\rho^{\alpha+ke}+\rho^\beta$ per $n$ dividitur, omnes ordinis termini erunt $\equiv(0)$ i.\,e. $\equiv1$. Q.\,E.\,D.

\emph{Annot.} Demonstratio haec simul methodum facillimam ostendit productum evolvendi. Aliam infra dabimus, quae hac quidem praerogativa caret, sed ob simplicitatem non contemnenda videtur.

\Article{248}
Periodos omnes minores, quae periodum maiorem constituunt, periodorum systema nominamus. Ita periodi
\[
(ef\ast\alpha),\quad(ef\ast f+\alpha),\quad(ef\ast2f+\alpha)\ \dotsfour\ (ef\ast(e-1)f+\alpha)
\]
e quibus componitur periodus $(f\ast\alpha)$, hoc nomine designabuntur. \emph{Rite ordinatum} erit, si numeri post signum $\ast$ positi, ut hic $\alpha$, $f+\alpha$, $2f+\alpha$, secundum seriem arithmeticam (cuius differentia est $f$) progrediantur; \emph{similia} denique erunt systemata, si tam minores quam maiores periodi sint similes.

\textsc{Theorema.} \emph{Si periodi systematum duorum similium rite ordinatorum invicem multiplicentur, prima scilicet in primam, secunda in secundam, tertia in tertiam etc., summa omnium productorum e periodis maiori similibus et numeris datis componi potest.}

\emph{Demonstr.} Sint systemata
\[
\begin{gathered}
(ef\ast\alpha),\quad(ef\ast\alpha+f),\quad(ef\ast\alpha+2f)\ldots\\
(ef\ast\beta),\quad(ef\ast\beta+f),\quad(ef\ast\beta+2f)\ldots
\end{gathered}
\]

\SourceLeaf{217}
Producta e singulis periodis systematis prioris in periodos respondentes posterioris constabunt (§.\ praec.) e numeris integris et periodis similibus. Sed parvula attentio ad genesin harum periodorum docebit, si $(ef\ast\alpha)\times(ef\ast\beta)$ constet ex numero integro $N$ et periodis $(ef\ast A)$, $(ef\ast B)$, $(ef\ast C)$ etc. tum constare producta
\[
\begin{aligned}
&(ef\ast\alpha+f)\times(ef\ast\beta+f)\text{ ex }N\text{ et perr.}\\
&\hspace{2em}(ef\ast A+f),\quad(ef\ast B+f),\quad(ef\ast C+f)\text{ etc.}\\
&(ef\ast\alpha+2f)\times(ef\ast\beta+2f)\text{ ex }N\text{ et perr.}\\
&\hspace{2em}(ef\ast A+2f),\quad(ef\ast B+2f),\quad(ef\ast C+2f)\text{ etc.}
\end{aligned}
\]
et generaliter
\[
\begin{aligned}
&(ef\ast\alpha+\mu f)\times(ef\ast\beta+\mu f)\text{ ex }N\text{ et perr.}\\
&\hspace{2em}(ef\ast A+\mu f),\quad(ef\ast B+\mu f),\quad(ef\ast C+\mu f)\text{ etc.}
\end{aligned}
\]
Unde sponte patet, omnium periodorum summam fore
\[
eN+(f\ast A)+(f\ast B)+(f\ast C)\text{ etc.}\qquad\text{Q.\,E.\,D.}
\]
Etiam haec demonstratio methodum suppeditat summam illam inveniendi.

\Article{249}
Facile est hoc theorema generalius adhuc reddere. Scilicet si habeantur quotcunque systemata rite ordinata similia fiantque producta ex omnibus periodis primis, secundis etc., omnium horum productorum summam constare e numeris et periodis maioribus. Si omnia haec systemata aequalia assumantur, summa potestatum quarumcunque omnium periodorum constabit e numeris et periodis maiori similibus. Iam hinc patescit, quorsum haec tendant. Sit $\lambda=efgh\ldots$; discerpantur omnes radices primae in $e$ periodos $A,A',A''$ etc., quaevis harum iterum in $f$: $B,B',B''$ etc., harum singulae in $g$: $C,C',C''$ etc. Iam omnium periodorum summa datur, est scilicet $\equiv-1$. Sed secundum ea, quae modo diximus, dabitur etiam
\[
\begin{gathered}
(A)^2+(A')^2+(A'')^2+(A''')^2+\text{etc.}\\
(A)^3+(A')^3+(A'')^3+(A''')^3+\text{etc.}\\
\text{etc. etc.}
\end{gathered}
\]

Hinc e §.\ 244 congruentia gradus $e^{\mathrm{ti}}$ inveniri poterit, cuius radices sint $A,A',A''$ etc. Iam his tamquam cognitis suppositis, quaevis periodus discerpatur in minores
\SourceLeaf{218}
\[
\begin{array}{rclll}
A&\text{in}&B,&B',&B''\ldots\\
A'&\text{in}&B^{(n)},&B^{(n+1)},&B^{(n+2)}\ldots\\
&\text{etc.}
\end{array}
\]
Datur ergo $B+B'+B''+\ldots\equiv A$. Sed constat
\[
\begin{gathered}
(B)^2+(B')^2+(B'')^2+\ldots\\
(B)^3+(B')^3+(B'')^3+\ldots\\
\text{etc.}
\end{gathered}
\]
ex unitatibus et periodis $A,A',A''$ etc. Quare $B,B',B''$ etc. dabuntur per congruentiam gradus $f^{\mathrm{ti}}$, ex qua inveniri possunt; similique modo periodi, ex quibus constant $A',A''$ etc., poterunt determinari. Quisquis autem hinc videbit, prorsus simili methodo quamvis periodum in minores subdividi posse, donec ad radices ipsas perveniatur.

\Article{250}
Sed in harum regularum applicatione difficultas occurrit, quam dimovere debemus. Quoniam scilicet quaevis congruentia plures radices habeat, quod cuique signum tribuendum sit, ut ab invicem rite dignosci possint, est videndum. Quoniam periodorum designatio a numeris $r,\rho$ pendet, qui ad libitum assumi possunt, necessario etiam designationi aliquid arbitrarii inhaerere debet. Numerus quidem $\rho$ iam ab initio est stabiliendus. Methodi nostrae indoles in eo potissimum consistit, ut ex periodis maioribus periodos minores deducamus. Sed hoc sine debito periodorum ordine, quem per \emph{signa} assecuti sumus, fieri nequit. Quare eo nitendum est, ut omnes periodi, quamprimum sunt inventae, signis suis distinguantur.

Sit periodus $A$ designata per $(e\ast\alpha)$ atque in $f$ periodos $B,C,D$ etc. discerpta, quas designare oportet. Patet quamvis in tali forma fore contentam $(ef\ast ke+\alpha)$; sed dico, pro aliqua earum $B$ numerum $k$ ad libitum assumi et inde ceterarum collocationem derivari posse.

Sit $R$ radix aliqua primitiva congr. $x^n\equiv1$ constetque $B$ e terminis $R^\mu+R^\nu+\text{etc.}$, sit $\frac1\mu\rho^{ke+\alpha}\equiv\frac1\upsilon\;(\mathrm{mod.}\,n)$ et quoniam valor ipsius $r$ est arbitrarius (si modo $A$ nanciscatur signum $(e\ast\alpha)$, quod sponte fieri manifestum est), ponatur $r\equiv R^\upsilon\;(\mathrm{mod.}\,p)$; quare terminus primus ipsius $B$ erit $r^{\rho^{ke+\alpha}}$ et $B$ per
\SourceLeaf{219}
$(ef\ast ke+\alpha)$ designare licet. Si loco ipsius $R^\mu$ terminum $R^\nu$ consideravissemus, alium ipsius $r$ valorem nacti essemus; sed sine negotio perspicitur, pro quacunque radice $\rho$, radicem $r$, $\frac\lambda{ef}$ valores diversos habere posse.

\Article{251}
Iam quomodo ex designatione unius periodi ceterae signis suis distinguantur, videamus. Ad hunc vero finem aliam methodum quaerere oportet reliquas periodos inveniendi; namque quatenus reliquae ut ipsa $A$ radices alicuius congruentiae sunt, nullus in illis ordo cernitur. Ponamus ipsum $A$ ita esse designatum $(ef\ast0)$, ex praecc. sequitur, fore
\[
\begin{aligned}
A^2&\text{ formae }M+N(ef\ast0)+O(ef\ast1)+P(ef\ast2)+\ldots\\
A^3&\text{ formae }M'+N'(ef\ast0)+O'(ef\ast1)+\ldots\\
&\text{etc.}\\
A^{ef-1}&\text{ formae }M^\ast+N^\ast(ef\ast0)+O^\ast(ef\ast1)+\ldots
\end{aligned}
\]
His accedit congruentia
\[
(ef\ast0)+(ef\ast1)+\ldots+(ef\ast ef-1)\equiv-1
\]
Habentur itaque $cf-1$ congruentiae lineares totidemque quantitates incognitae, quae igitur per eliminationem determinari possunt.

\emph{Annot.} Casus occurrere potest, quo quantitates incognitae per huiusmodi expressiones dantur $\frac{V}{Wp}$; quomodo vero huic difficultati remedium afferri possit, infra docebimus. Hic, quoniam hic casus perraro occurrere potest, ei immorari nolumus.

\Article{252}
Haec in genere de solutione congruentiarum purarum sufficiant. Passim infra multa adhuc de ipsis dicentur; praesertim multa ex solutione aequationum purarum huc trahi possunt, quae loco suo annotare non negligemus. Exemplum adhuc apponimus, quo cum praeceptis collato, omnia minus peritis clariora fient.

Sit $n=31$, $p=311$, sive investigandae sunt radices congruentiae $x^{31}-1\equiv0\;(\mathrm{mod.}\,311)$. Statim radix primitiva congruentiae $y^{30}-1\equiv\;(\mathrm{mod.}\,31)$ est quaerenda, qualis est $y\equiv3$. Ponamus itaque $\rho\equiv3$ et omnes congruentiae propositae radices primitivas primum in $5$ periodos discerpamus, scilicet
\SourceLeaf{220}
\[
\begin{array}{rl}
(5\ast0)\ \dotsfive&(1)+(26)+(25)+(30)+(5)+(6)\\
(5\ast1)\ \dotsfive&(3)+(16)+(13)+(28)+(15)+(18)\\
(5\ast2)\ \dotsfive&(9)+(17)+(8)+(22)+(14)+(23)\\
(5\ast3)\ \dotsfive&(27)+(20)+(24)+(4)+(11)+(7)\\
(5\ast4)\ \dotsfive&(19)+(29)+(10)+(12)+(2)+(21)
\end{array}
\]

Per calculos requisitos invenietur summa periodd. $\equiv-1$, quadrat. $\equiv25$, cub. $\equiv26$, biquad. $\equiv249$, pott. quintt. $\equiv564$.

Quare periodi erunt radices congruentiae
\[
x^5+x^4-12x^3-21x^2+x+5\equiv0
\]

Porro autem invenitur
\[
\begin{aligned}
(5\ast0)^2&\equiv6+2(5\ast0)+2(5\ast3)+(5\ast4)\\
(5\ast0)^3&\equiv12+15(5\ast0)+4(5\ast1)+3(5\ast2)+6(5\ast3)+6(5\ast4)\\
(5\ast0)^4&\equiv90+60(5\ast0)+28(5\ast1)+26(5\ast2)+49(5\ast3)+38(5\ast4)
\end{aligned}
\]
et hinc per eliminationem
\[
\begin{aligned}
5(5\ast1)&\equiv3(5\ast0)^4-(5\ast0)^3-33(5\ast0)^2-24(5\ast0)+15\\
5(5\ast2)&\equiv-2(5\ast0)^4-(5\ast0)^3+22(5\ast0)^2+31(5\ast0)\\
5(5\ast3)&\equiv(5\ast0)^4-2(5\ast0)^3\\
5(5\ast4)&\equiv-2(5\ast0)^4+4(5\ast0)^3
\end{aligned}
\]

Congruentiae vero inventae una radix est $\equiv17$; quare si ponatur $(5\ast0)\equiv17$, erit $(5\ast1)\equiv183$, $(5\ast2)\equiv263$, $(5\ast3)\equiv91$, $(5\ast4)\equiv67$.

Iam periodi inventae iterum discerpantur singulae in ternas; scilicet
\[
\begin{array}{rl}
(5\ast0)&\text{in }(15\ast0),\ (15\ast5),\ (15\ast10)\\
&\text{sive in }(1)+(30),\ (26)+(5),\ (25)+(6)\\
(5\ast1)&\text{in }(15\ast1),\ (15\ast6),\ (15\ast11)\\
&\text{sive in }(3)+(28),\ (16)+(15),\ (13)+(18)\\
&\text{etc.}\qquad\text{etc.}
\end{array}
\]

Ponatur periodos, in quas discerpta est
\[
\begin{array}{rll}
(5\ast0)&\text{esse radices congr.}&x^3+Ax^2+Bx+C\equiv0\\
(5\ast1)&&x^3+A'x^2+B'x+C'\equiv0\\
(5\ast2)&&x^3+A''x^2+B''x+C''\equiv0\\
&\text{etc.}
\end{array}
\]

\SourceLeaf{221}
eritque
\[
\begin{array}{lll}
A\equiv-(5\ast0),&B\equiv(5\ast0)+(5\ast3),&C\equiv-2-(5\ast4)\\
A'\equiv-(5\ast1),&B'\equiv(5\ast1)+(5\ast4),&C'\equiv-2-(5\ast0)\\
\text{etc.}&\text{etc.}&\text{etc.}
\end{array}
\]

Quare
\[
\begin{array}{lll}
(15\ast0),\ (15\ast5),\ (15\ast10)&\begin{gathered}\text{erunt radices}\\\text{congr.}\end{gathered}&x^3-17x^2+108x-60\equiv0\\
(15\ast1),\ (15\ast6),\ (15\ast11)&&\\
(15\ast2),\ (15\ast7),\ (15\ast12)&&\\
(15\ast3),\ (15\ast8),\ (15\ast13)&&\\
(15\ast4),\ (15\ast9),\ (15\ast14)&&
\end{array}
\]

Hic autem habetur
\[
\begin{gathered}
(15\ast0)^3-3(15\ast0)\equiv(15\ast1)\\
(15\ast1)^3-3(15\ast1)\equiv(15\ast2)\\
\text{etc.}
\end{gathered}
\]

Unde si una radicum primae congruentiae, $10$, ponatur $(15\ast0)$ habetur
\[
\begin{array}{rlrlrl}
(15\ast0)\equiv&10&(15\ast5)\equiv&& (15\ast10)\equiv&\\
(15\ast1)\equiv&37&(15\ast6)\equiv&& (15\ast11)\equiv&\\
(15\ast2)\equiv&-151&(15\ast7)\equiv&& (15\ast12)\equiv&\\
(15\ast3)\equiv&-39&(15\ast8)\equiv&& (15\ast13)\equiv&\\
(15\ast4)\equiv&-112&(15\ast9)\equiv&& (15\ast14)\equiv&
\end{array}
\]

Tandem harum singularum periodorum capiantur termini constituentes eruntque
\[
\begin{array}{lll}
(1),\ (30)&\text{radices congr.}&x^2-(15\ast0)x+1\equiv0\\
(3),\ (28)&&x^2-(15\ast1)x+1\equiv0\\
&\text{etc.}
\end{array}
\]

Prima congruentiae radices sunt $126$ et $195$, quae igitur erunt radices primitivae congruentiae $x^{31}\equiv1$ et ex his reliquae sine negotio deduci possunt.
\BookRule{.28\linewidth}
